\documentclass[10pt,a4paper]{amsart}
\usepackage{amsmath,amssymb,amsthm,graphicx}
% Retained minimal preamble for the Stacks Project chapter more-algebra.tex (More on Algebra).
% The chapter file pinned-original-source.tex carries only \input{preamble} on its first line, so no
% preamble text lives inside the chapter itself.  The Stacks Project uses one single shared preamble for
% every chapter; the snapshot stages it identically as sources/stacks-chow/preamble.tex and
% sources/stacks-etale-cohomology/preamble.tex.  The block below is copied mechanically, line for line,
% from sources/stacks-chow/preamble.tex lines 14-18 (the Xy-pic setup) and 154-217 (the native theorem
% declarations together with the mathematical macros).  Nothing was added, removed or redefined; the
% wrapper class and the amsmath/amssymb/amsthm/graphicx base packages are supplied by the assembler.
% Chapter: more-algebra.tex -- More on Algebra.
% For commutative diagrams we use Xy-pic
\usepackage[all]{xy}

% We use 2cell for 2-commutative diagrams.
\xyoption{2cell}

% Theorem environments.
%
\theoremstyle{plain}
\newtheorem{theorem}[subsection]{Theorem}
\newtheorem{proposition}[subsection]{Proposition}
\newtheorem{lemma}[subsection]{Lemma}

\theoremstyle{definition}
\newtheorem{definition}[subsection]{Definition}
\newtheorem{example}[subsection]{Example}
\newtheorem{exercise}[subsection]{Exercise}
\newtheorem{situation}[subsection]{Situation}

\theoremstyle{remark}
\newtheorem{remark}[subsection]{Remark}
\newtheorem{remarks}[subsection]{Remarks}

\numberwithin{equation}{subsection}

% Macros
%
\def\lim{\mathop{\mathrm{lim}}\nolimits}
\def\colim{\mathop{\mathrm{colim}}\nolimits}
\def\Spec{\mathop{\mathrm{Spec}}}
\def\Hom{\mathop{\mathrm{Hom}}\nolimits}
\def\Ext{\mathop{\mathrm{Ext}}\nolimits}
\def\SheafHom{\mathop{\mathcal{H}\!\mathit{om}}\nolimits}
\def\SheafExt{\mathop{\mathcal{E}\!\mathit{xt}}\nolimits}
\def\Sch{\mathit{Sch}}
\def\Mor{\mathop{\mathrm{Mor}}\nolimits}
\def\Ob{\mathop{\mathrm{Ob}}\nolimits}
\def\Sh{\mathop{\mathit{Sh}}\nolimits}
\def\NL{\mathop{N\!L}\nolimits}
\def\CH{\mathop{\mathrm{CH}}\nolimits}
\def\proetale{{pro\text{-}\acute{e}tale}}
\def\etale{{\acute{e}tale}}
\def\QCoh{\mathit{QCoh}}
\def\Ker{\mathop{\mathrm{Ker}}}
\def\Im{\mathop{\mathrm{Im}}}
\def\Coker{\mathop{\mathrm{Coker}}}
\def\Coim{\mathop{\mathrm{Coim}}}

% Boxtimes
%
\DeclareMathSymbol{\boxtimes}{\mathbin}{AMSa}{"02}

%
% Macros for moduli stacks/spaces
%
\def\QCohstack{\mathcal{QC}\!\mathit{oh}}
\def\Cohstack{\mathcal{C}\!\mathit{oh}}
\def\Spacesstack{\mathcal{S}\!\mathit{paces}}
\def\Quotfunctor{\mathrm{Quot}}
\def\Hilbfunctor{\mathrm{Hilb}}
\def\Curvesstack{\mathcal{C}\!\mathit{urves}}
\def\Polarizedstack{\mathcal{P}\!\mathit{olarized}}
\def\Complexesstack{\mathcal{C}\!\mathit{omplexes}}
% \Pic is the operator that assigns to X its picard group, usage \Pic(X)
% \Picardstack_{X/B} denotes the Picard stack of X over B
% \Picardfunctor_{X/B} denotes the Picard functor of X over B
\def\Pic{\mathop{\mathrm{Pic}}\nolimits}
\def\Picardstack{\mathcal{P}\!\mathit{ic}}
\def\Picardfunctor{\mathrm{Pic}}
\def\Deformationcategory{\mathcal{D}\!\mathit{ef}}

% --- appended 2026-10-12 (Chief): the project own definitions of the reference/slogan/history
% display environments, copied mechanically from the shared Stacks preamble lines 8-12
% (\usepackage{verbatim} plus the three \newenvironment lines, which the chapter preamble keeps
% verbatim from the project preamble). Nothing added, removed or redefined: these environments
% carry no mathematical content and the project itself typesets them as comments. Without these
% three lines a counted statement whose own text uses them fails with Environment slogan undefined.
\usepackage{verbatim}
\newenvironment{reference}{\comment}{\endcomment}
\newenvironment{slogan}{\comment}{\endcomment}
\newenvironment{history}{\comment}{\endcomment}
\usepackage{cleveref}
\crefname{originallemma}{original lemma}{original lemmas}
\Crefname{originallemma}{Original lemma}{Original lemmas}
\makeatletter
\newlabel{algebra-lemma-nagata-pth-roots}{{09E2}{source}{}{sourcepointer-1}{}}
\newlabel{algebra-lemma-nagata-pth-roots@cref}{{[originallemma][0][]09E2}{[1][0][]source}{}{}{}}
\makeatother
\crefname{originallemma}{original lemma}{original lemmas}
\Crefname{originallemma}{Original lemma}{Original lemmas}
\makeatletter
\newlabel{lemma-inequality}{{32936}{source}{}{sourcepointer-2}{}}
\newlabel{lemma-inequality@cref}{{[originallemma][0][]32936}{[1][0][]source}{}{}{}}
\makeatother
\crefname{originallemma}{original lemma}{original lemmas}
\Crefname{originallemma}{Original lemma}{Original lemmas}
\makeatletter
\newlabel{lemma-make-degree-q-extension}{{34641}{source}{}{sourcepointer-3}{}}
\newlabel{lemma-make-degree-q-extension@cref}{{[originallemma][0][]34641}{[1][0][]source}{}{}{}}
\makeatother
\crefname{originallemma}{original lemma}{original lemmas}
\Crefname{originallemma}{Original lemma}{Original lemmas}
\makeatletter
\newlabel{lemma-pre-purely-inseparable-case}{{34663}{source}{}{sourcepointer-4}{}}
\newlabel{lemma-pre-purely-inseparable-case@cref}{{[originallemma][0][]34663}{[1][0][]source}{}{}{}}
\makeatother
\crefname{originallemma}{original lemma}{original lemmas}
\Crefname{originallemma}{Original lemma}{Original lemmas}
\makeatletter
\newlabel{lemma-ramification-index-a-power-of-p}{{32965}{source}{}{sourcepointer-5}{}}
\newlabel{lemma-ramification-index-a-power-of-p@cref}{{[originallemma][0][]32965}{[1][0][]source}{}{}{}}
\makeatother
\crefname{originallemma}{original lemma}{original lemmas}
\Crefname{originallemma}{Original lemma}{Original lemmas}
\makeatletter
\newlabel{lemma-weakly-unramified-goes-up-along-totally-ramified}{{34490}{source}{}{sourcepointer-6}{}}
\newlabel{lemma-weakly-unramified-goes-up-along-totally-ramified@cref}{{[originallemma][0][]34490}{[1][0][]source}{}{}{}}
\makeatother
\title{For a degree p purely inseparable step over a weakly unramified extension one of three explicitly described alternatives holds}
\author{The Stacks Project Authors}
\date{stacks-project:more-algebra.tex:local-snapshot-20261012; GNU FDL 1.2 or later}
\begin{document}
\maketitle
\setcounter{subsection}{1}
\noindent\textit{Editorial note.} All delivered mathematical text is the author's own, copied from the pinned Stacks Project chapter file more-algebra.tex; no statement or proof step was rewritten, reordered or removed.\par\medskip
\noindent\textit{Editorial note.} The chapter file pinned-original-source.tex carries only \textbackslash{}input\{preamble\} on its first line, so no preamble text lives inside the chapter. The retained file sources/stacks-more-algebra/retained-preamble.tex is a mechanical, line-for-line copy of the project's single shared preamble as staged in the snapshot; nothing was added, removed or redefined and no package was added.\par\medskip
\noindent\textit{Editorial note.} The source URL is the source-level Stacks Project URL, as required when more than one card is produced from one source; the per-claim tag is carried by claim\_key.\par\medskip
\noindent\textit{Editorial note.} The author writes '(computation omitted)' inside the counted proof span, reproduced verbatim. The six references inside the counted proof span are recorded as non-load-bearing navigation pointers.\par\medskip
\noindent\textit{Allowed standard prerequisites (editorial).}\begin{itemize}
\item Ramification indices of a tower of discrete valuation rings and their multiplicativity, with the ramification index of a purely inseparable degree p step being 1 or p. used at proof lines 34698-34701. Reduces the case analysis to e = p and gives the agreement of the residue fields.
\item Weakly unramified extensions and their behaviour under a totally ramified base change. the hypothesis at statement line 34682 and used at proof lines 34718-34719. Is what is preserved when passing to K\_1 and B\_1, and is one of the conclusions in the third alternative.
\item The degree-q extension construction with prescribed uniformizer relation. invoked at proof lines 34712-34716. Produces the auxiliary K\_1 and the uniformizer pi' used in the computation.
\item Nagata rings and the descent of p-th roots from the completion. the hypothesis at statement line 34681 and used at proof lines 34740-34742. Is what settles the second case.
\end{itemize}\medskip
\ifdefined\hypertarget\hypertarget{sourcepointer-1}{}\fi
\noindent\textit{Source reference.} Original lemma 09E2: focused-seeds/AGP/sources/stacks-snapshot-20261011/tags.txt, source lines 11885--11885. Cross-chapter navigation pointer into the Stacks Project; the standard fact it supplies is applied inside the counted proof span and the argument resting on it stays there.\par\medskip
\ifdefined\hypertarget\hypertarget{sourcepointer-2}{}\fi
\noindent\textit{Source reference.} Original lemma 32936: sources/stacks-more-algebra/pinned-original-source.tex, source lines 32936--32936. Same-chapter navigation pointer; the statement it names is standard for the section and the argument that consumes it stays inside the counted proof span.\par\medskip
\ifdefined\hypertarget\hypertarget{sourcepointer-3}{}\fi
\noindent\textit{Source reference.} Original lemma 34641: sources/stacks-more-algebra/pinned-original-source.tex, source lines 34641--34641. Same-chapter navigation pointer; the statement it names is standard for the section and the argument that consumes it stays inside the counted proof span.\par\medskip
\ifdefined\hypertarget\hypertarget{sourcepointer-4}{}\fi
\noindent\textit{Source reference.} Original lemma 34663: sources/stacks-more-algebra/pinned-original-source.tex, source lines 34663--34663. Same-chapter navigation pointer; the statement it names is standard for the section and the argument that consumes it stays inside the counted proof span.\par\medskip
\ifdefined\hypertarget\hypertarget{sourcepointer-5}{}\fi
\noindent\textit{Source reference.} Original lemma 32965: sources/stacks-more-algebra/pinned-original-source.tex, source lines 32965--32965. Same-chapter navigation pointer; the statement it names is standard for the section and the argument that consumes it stays inside the counted proof span.\par\medskip
\ifdefined\hypertarget\hypertarget{sourcepointer-6}{}\fi
\noindent\textit{Source reference.} Original lemma 34490: sources/stacks-more-algebra/pinned-original-source.tex, source lines 34490--34490. Same-chapter navigation pointer; the statement it names is standard for the section and the argument that consumes it stays inside the counted proof span.\par\medskip
\begin{lemma}
\label{lemma-purely-inseparable-case}
Let $A \subset B \subset C$ be extensions of discrete valuation rings
with fractions fields $K \subset L \subset M$. Let $\pi \in A$ be a
uniformizer. Assume
\begin{enumerate}
\item $B$ is a Nagata ring,
\item $A \subset B$ is weakly unramified,
\item $M$ is a degree $p$ purely inseparable extension of $L$.
\end{enumerate}
Then either
\begin{enumerate}
\item $A \to C$ is weakly unramified, or
\item $C = B[\pi^{1/p}]$, or
\item there exists a degree $p$ separable extension $K_1/K$
totally ramified with respect to $A$
such that $L_1 = L \otimes_K K_1$ and $M_1 = M \otimes_K K_1$
are fields and the maps of integral closures $A_1 \to B_1 \to C_1$
are weakly unramified extensions of discrete valuation rings.
\end{enumerate}
\end{lemma}

\begin{proof}
Let $e$ be the ramification index of $C$ over $B$. If $e = 1$, then we are
done. If not, then $e = p$ by Lemmas \ref{lemma-inequality} and
\ref{lemma-ramification-index-a-power-of-p}.
This in turn implies that the residue fields of $B$ and $C$ agree.
Choose a uniformizer $\pi_C$ of $C$.
Write $\pi_C^p = u \pi$ for some unit $u$ of $C$.
Since $\pi_C^p \in L$, we see that $u \in B^*$. Also $M = L[\pi_C]$.

\medskip\noindent
Suppose there exists an integer $m \geq 0$ such that
$$
u = \sum\nolimits_{0 \leq i < m} b_i^p \pi^i + b \pi^m
$$
with $b_i \in B$ and with $b \in B$ an element whose image in $\kappa_B$
is not a $p$th power. Choose an extension $K_1/K$ as in
Lemma \ref{lemma-make-degree-q-extension}
with $n = m + 2$ and denote $\pi'$ the uniformizer
of the integral closure $A_1$ of $A$ in $K_1$ such that
$\pi = (\pi')^p + (\pi')^{np} a$ for some $a \in A_1$.
Let $B_1$ be the integral closure of $B$ in $L \otimes_K K_1$.
Observe that $A_1 \to B_1$ is weakly unramified by
Lemma \ref{lemma-weakly-unramified-goes-up-along-totally-ramified}.
In $B_1$ we have
$$
u \pi =
\left(\sum\nolimits_{0 \leq i < m} b_i (\pi')^{i + 1}\right)^p +
b (\pi')^{(m + 1)p} + (\pi')^{np} b_1
$$
for some $b_1 \in B_1$ (computation omitted).
We conclude that $M_1$ is obtained from
$L_1$ by adjoining a $p$th root of
$$
b + (\pi')^{n - m - 1} b_1
$$
Since the residue field of $B_1$ equals the residue field of $B$
we see from Lemma \ref{lemma-pre-purely-inseparable-case}
that $M_1/L_1$ has degree $p$ and
the integral closure $C_1$ of $B_1$ is weakly unramified over $B_1$.
Thus we conclude in this case.

\medskip\noindent
If there does not exist an integer $m$ as in the preceding paragraph,
then $u$ is a $p$th power in the $\pi$-adic completion of $B_1$.
Since $B$ is Nagata, this means that $u$ is a $p$th power in $B_1$
by Algebra, Lemma \ref{algebra-lemma-nagata-pth-roots}.
Whence the second case of the statement of the lemma holds.
\end{proof}

\end{document}
